\documentclass[12pt]{article}
\usepackage{amsmath}
\usepackage{amssymb}
\usepackage{geometry}
\usepackage{setspace}
\geometry{margin=1in}

\title{\textbf{CSE400 – Fundamentals of Probability in Computing}\\
\textbf{Lecture 10: Randomized Min-Cut Algorithm}}
\author{Dhaval Patel, PhD\\
Associate Professor\\
SEAS – Ahmedabad University, Ahmedabad, Gujarat, India\\[1em]
\textbf{Nishka Shah – AU2440168}}
\date{February 5, 2026}

\begin{document}
\maketitle

\section*{Outline}

\begin{itemize}
    \item \textbf{Min-Cut Problem}
    \begin{itemize}
        \item Why use min-cut?
        \item What is min-cut?
        \item Successful min-cut run
        \item Unsuccessful min-cut run
        \item Max-Flow Min-Cut Theorem
    \end{itemize}
    \item \textbf{Deterministic Min-Cut Algorithm}
    \begin{itemize}
        \item Stoer–Wagner minimum cut algorithm
        \item Pseudocode
    \end{itemize}
    \item \textbf{Randomized Min-Cut Algorithm}
    \begin{itemize}
        \item Why randomized algorithm?
        \item Karger’s Randomized Algorithm
        \item Pseudocode
        \item Comparison: Deterministic vs Randomized
        \item Theorem for min-cut set
        \item Python Simulation
    \end{itemize}
\end{itemize}

\section*{Min-Cut Problem}

\subsection*{Why use min-cut?}

We use the min-cut algorithm in various applications to solve problems related to \textbf{network connectivity, reliability, and optimization}.

\textbf{Network Design:}  
Min cut helps in improving the efficiency of communication and optimizing network flow. The algorithm is used in network design to find the minimum capacity cut.

\textbf{Communication Networks:}  
For understanding the vulnerability of networks to failures, minimum cut can be useful. It helps in building robust and fault-tolerant communication networks.

\textbf{VLSI Design:}  
In Very Large Scale Integration (VLSI) design, the algorithm is useful for partitioning circuits into smaller components leading to reduced interconnectivity complexity.

Refer to \textbf{Section 1.5. Application: A Randomized Min-Cut Algorithm} in \textit{Probability and Computing (2nd Edition)}.

\subsection*{What is Min-Cut?}

A \textbf{cut-set} in a graph is a set of edges whose removal breaks the graph into two or more connected components.

Given a graph  
\[
G = (V, E)
\]
with \( n \) vertices, the \textbf{minimum cut (min-cut)} problem is to find a \textbf{minimum cardinality cut-set} in \( G \).

Min-cut algorithms such as \textbf{Karger’s algorithm} are random and can be sensitive to the initial choice of edges. If the algorithm contracts critical edges early, it may fail to find the true minimum cut.

The main operation is \textbf{edge contraction}.

\textbf{Edge contraction:}  
Contracting an edge \( (u,v) \):
\begin{itemize}
    \item Merge vertices \( u \) and \( v \) into one vertex
    \item Remove all edges between \( u \) and \( v \)
    \item Retain all other edges
    \item Parallel edges may occur
    \item Self-loops are removed
\end{itemize}

\subsection*{Successful Min-Cut Run}

A successful min-cut run refers to the success in the outcome of an algorithm designed to find the minimum cut in a graph.

(Figure: Successful Run)

\subsection*{Unsuccessful Min-Cut Run}

An unsuccessful min-cut run refers to an iteration where the algorithm fails to correctly identify the minimum cut of a given graph.

(Figure: Unsuccessful Run)

\section*{Max-Flow Min-Cut Theorem}

The theorem states:

\begin{quote}
\textbf{In a flow network, the maximum amount of flow from the source to the sink is equal to the total weight of the edges in a minimum cut.}
\end{quote}

\textbf{Capacity of a cut:}  
Sum of capacities of edges oriented from vertices
\[
\in X \text{ to } \in Y
\]

\textbf{Minimum cut:}  
The cut with the smallest possible capacity.

\textbf{Minimum cut capacity:}  
Capacity of the minimum cut.

\textbf{Maximum flow:}  
Largest possible flow from source \( S \) to sink \( T \).

\section*{Deterministic Min-Cut Algorithm}

\subsection*{Stoer–Wagner Min-Cut Algorithm}

Let \( s \) and \( t \) be two vertices of graph \( G \).  
Let \( G/\{s,t\} \) be the graph obtained by merging \( s \) and \( t \).

Then a minimum cut of \( G \) is the smaller of:
\begin{itemize}
    \item A minimum \( s\text{-}t \) cut of \( G \)
    \item A minimum cut of \( G/\{s,t\} \)
\end{itemize}

The theorem holds because:
\begin{itemize}
    \item Either a minimum cut separates \( s \) and \( t \)
    \item Or it does not, and contraction preserves the minimum cut
\end{itemize}

\subsection*{Pseudocode}

\textbf{Algorithm 1: MinimumCutPhase(G, a)}

\begin{verbatim}
A ← {a}
while A ≠ V do
    add to A the most tightly connected vertex
return the cut weight as the “cut of the phase”
\end{verbatim}

\textbf{Algorithm 2: MinimumCut(G)}

\begin{verbatim}
while |V| ≥ 1 do
    choose any a from V
    MinimumCutPhase(G, a)
    if cut-of-the-phase is lighter than current minimum cut then
        store cut-of-the-phase
    shrink G by merging the two vertices added last
return the minimum cut
\end{verbatim}

\section*{Randomized Min-Cut Algorithm}

\subsection*{Why Randomized Algorithm?}

\begin{itemize}
    \item Efficiency
    \item Parallelization
    \item Approximation Guarantees
    \item Avoidance of Worst-Case Instances
    \item Heuristic Nature
    \item Robustness
\end{itemize}

\subsection*{Karger’s Randomized Algorithm}

(Randomized Min-Cut Algorithm: Karger’s Randomized Algorithm)

\section*{Comparison: Deterministic vs Randomized Min-Cut}

The choice depends on the problem.

\textbf{Deterministic Min-Cut:}
\begin{itemize}
    \item Guarantees exact minimum cut
    \item Higher time complexity
    \item Stoer-Wagner algorithm complexity:
    \[
    O(V \cdot E + V^2 \log V)
    \]
\end{itemize}

\textbf{Randomized Min-Cut:}
\begin{itemize}
    \item Gives approximate min-cut with high probability
    \item Karger’s algorithm complexity:
    \[
    O(V^2)
    \]
\end{itemize}

\section*{Theorem for Min-Cut Set}

The algorithm outputs a min-cut set with probability at least:

\[
\frac{2}{n(n-1)}
\]

\section*{Python Simulation}

\begin{itemize}
    \item Open Campuswire post for Lecture 10
    \item Download the provided \texttt{.ipynb} file
\end{itemize}

\section*{Thank You}

\textbf{Dhaval Patel, PhD}\\
CSE400

\end{document}
